%! Inference and Experiments — Unit 3, Lesson 3.5 — Warm-Up (Answer Key)
% Mirrors the blank exactly apart from the package swap, the pageheader, and the answers.
% NO teachernote here — teacher prose lives in the lesson plan.
% ONE PAGE, blank and key. Three items; do not add a fourth.
% GRADUAL RELEASE: the warm-up activates prior knowledge before the guided notes.
%   Item 1 — spiral to Lessons 3.2 and 3.4: random selection or random assignment, and what each
%            one is for. Seeds the Scope of Inference row (selection -> generalize,
%            assignment -> cause).
%   Item 2 — spiral to Lesson 3.1: why an observational study cannot show cause (confounding).
%            Seeds the crux row (Association or Cause?).
%   Item 3 — the music experiment: could a 2.5-word difference come from the luck of the random
%            assignment alone? A vote, not a verdict — leave the room's answer on the board; the
%            simulation dotplot in notes row 1 decides it.
% Every display is given; nothing is sketched.
\documentclass[12pt]{article}
\usepackage{apstats-article}
\usepackage{apstats-key}   % pulls in -boxes; do NOT also load -boxes

\begin{document}

\pageheader{Unit 3, Lesson 3.5}{Warm-Up --- Answer Key}
% NAMESTRIP: no \namedateperiod here — the name row lives on the cover only.

\vspace{0.28in}

\begin{enumerate}[leftmargin=1.9em, itemsep=0.42in]

  \item \textbf{From Lessons 3.2 and 3.4.} Name the kind of randomness used in each: random
        \emph{selection} or random \emph{assignment}.

        \vspace{10pt}
        (a) 50 names are drawn at random from the school roster to survey. \hfill \ans{Random selection}

        \vspace{10pt}
        (b) A coin flip decides which of 30 volunteers get a new drug. \hfill \ans{Random assignment}

        \vspace{12pt}
        What is each kind of randomness \emph{for}?
        \vspace{12pt}

        \ansline{Selection makes the sample represent the population. Assignment makes}\\[18pt]\ansline{the groups alike, so the treatment is the only planned difference.}

  \item \textbf{From Lesson 3.1.} An observational study finds that students who eat breakfast
        tend to earn higher grades. Why can't we conclude that eating breakfast \emph{causes}
        higher grades?

        \vspace{12pt}
        \ansline{No one was assigned breakfast. A confounding variable, such as a steady}\\[18pt]\ansline{home routine, could lead to both breakfast and higher grades.}

  \item 24 volunteers were randomly assigned to study a word list in silence or with music, 12
        in each group. The silence group recalled \textbf{2.5 more words} on average. Could a
        difference this big happen just from the luck of \emph{which} volunteers landed in
        which group, even if music makes no difference at all? Explain your thinking.

        \vspace{12pt}
        \ansline{Possibly. By luck the random assignment might put stronger memorizers in}\\[18pt]\ansline{the silence group. We need to know how often chance alone gives 2.5 or}\\[18pt]\ansline{more --- the notes find out.}

\end{enumerate}

\end{document}
